% Every script this build ships a face for, wired up so that a note in one of them draws with no
% reader action at all.
%
% **Not the reader's file.** Yi writes this one into its export folder before each PDF export, with the
% hole below — written between double at-signs, where the faces are defined — filled in with one face
% per script, by this build's path or by a name the reader's settings file gave. It goes to pandoc with
% `--include-in-header`: added to, not instead of, whatever the reader's own settings file lists there.
% The reader's file is never edited, so a change of build cannot reach into it; and nothing here is
% meant to be edited by hand, either — a face of your own is a line in the settings file, not a change
% to this one.
%
% **A font the reader names wins**, and there are two ways to name one: `mainfont`, the document's own
% font — `\normalfont` below is that font, so it is what every Latin character and every block no face
% is defined for is set in — and, for one script at a time, the settings file's `yi-<script>-font`
% line (Han and kana answer to `CJKmainfont`, which is pandoc's own key). A script nobody names is
% drawn with the face this build ships, so a note in one of them needs no line at all.
%
% **Why classes at all.** XeTeX has no font fallback: a character outside the current font is
% reported missing and drawn as nothing. `ucharclasses` gives each Unicode block a character class
% and XeTeX an interchartoks hook for entering one, so entering a script's block can set that
% script's face. Only *entering* hooks are written: a "leaving" hook is written per class **pair**,
% so two classes that both write one each overwrite the other's entering hook (measured — that is
% how the first version lost Greek after Cyrillic).
%
% **Blocks, not groups.** `ucharclasses` has informal groups (`Latin`, `CJK`, `Arabics`, …) but no
% group at all for Hebrew, Thai or Armenian, and a hook written for a block the package did not load
% is a hook that silently does nothing. So every block is named.
\ifdefined\XeTeXinterchartoks
\usepackage{ucharclasses}
% Break lines between CJK characters, which TeX does not do on its own: without this a Han
% paragraph runs off the page as one line, measured both ways.
\XeTeXlinebreaklocale "zh"
\XeTeXlinebreakskip = 0pt plus 1pt
% One face per script, in the order the header draws them: this build's own face by file, or — where
% the settings file names one for that script — the reader's, by name. `export_fonts::write_header`
% writes this block, so the list of scripts and the settings file's keys for them are one fact, in
% that module, and the routing below is the other.
@@FACES@@
% The document's own text: `\normalfont` is the file's own font (or the engine's), here as
% everywhere, so a reader who names `mainfont` keeps it.
\setTransitionTo{BasicLatin}{\normalfont}
\setTransitionTo{LatinSupplement}{\normalfont}
\setTransitionTo{LatinExtendedA}{\normalfont}
\setTransitionTo{LatinExtendedB}{\normalfont}
\setTransitionTo{LatinExtendedC}{\normalfont}
\setTransitionTo{LatinExtendedD}{\normalfont}
\setTransitionTo{LatinExtendedE}{\normalfont}
\setTransitionTo{LatinExtendedF}{\normalfont}
\setTransitionTo{LatinExtendedG}{\normalfont}
\setTransitionTo{LatinExtendedAdditional}{\normalfont}
% Greek and Cyrillic: one face, the one that has both.
\setTransitionTo{GreekAndCoptic}{\yiGreek}
\setTransitionTo{GreekExtended}{\yiGreek}
\setTransitionTo{Coptic}{\yiGreek}
\setTransitionTo{CopticEpactNumbers}{\yiGreek}
\setTransitionTo{Cyrillic}{\yiCyrillic}
\setTransitionTo{CyrillicSupplement}{\yiCyrillic}
\setTransitionTo{CyrillicExtendedA}{\yiCyrillic}
\setTransitionTo{CyrillicExtendedB}{\yiCyrillic}
\setTransitionTo{CyrillicExtendedC}{\yiCyrillic}
\setTransitionTo{CyrillicExtendedD}{\yiCyrillic}
\setTransitionTo{Glagolitic}{\yiCyrillic}
\setTransitionTo{GlagoliticSupplement}{\yiCyrillic}
\setTransitionTo{Hebrew}{\yiHebrew}
% The block the Hebrew face is asked for beyond the Hebrew block itself: it is where the Hebrew
% presentation forms live (measured: the face holds `U+FB1D`–`U+FB4F`), and it is one block with Latin
% and Armenian ligature code points in it as well — a corner of the text few notes use, and the price of
% the block being one class rather than three.
\setTransitionTo{AlphabeticPresentationForms}{\yiHebrew}
\setTransitionTo{Arabic}{\yiArabic}
\setTransitionTo{ArabicSupplement}{\yiArabic}
\setTransitionTo{ArabicExtendedA}{\yiArabic}
\setTransitionTo{ArabicExtendedB}{\yiArabic}
\setTransitionTo{ArabicExtendedC}{\yiArabic}
\setTransitionTo{ArabicPresentationFormsA}{\yiArabic}
\setTransitionTo{ArabicPresentationFormsB}{\yiArabic}
\setTransitionTo{ArabicMathematicalAlphabeticSymbols}{\yiArabic}
\setTransitionTo{DevanagariPreMarks}{\yiDevanagari}
\setTransitionTo{DevanagariMarks}{\yiDevanagari}
\setTransitionTo{DevanagariPostMarks}{\yiDevanagari}
\setTransitionTo{DevanagariDanDa}{\yiDevanagari}
\setTransitionTo{DevanagariPostDanDa}{\yiDevanagari}
\setTransitionTo{DevanagariExtended}{\yiDevanagari}
\setTransitionTo{DevanagariExtendedA}{\yiDevanagari}
\setTransitionTo{Thai}{\yiThai}
\setTransitionTo{Armenian}{\yiArmenian}
\setTransitionTo{Georgian}{\yiGeorgian}
\setTransitionTo{GeorgianSupplement}{\yiGeorgian}
\setTransitionTo{GeorgianExtended}{\yiGeorgian}
% Hangul: this build's own face unless the file names one for it, which is half the reason a face is
% shipped at all — the fonts a platform has for Han very often have none of it (measured: the macOS
% `Songti SC` draws Han and kana and no Hangul, and a Korean note that named it lost every syllable).
% A `CJKmainfont` — Han's own line — does not reach here, and neither does a `mainfont`.
\setTransitionTo{HangulSyllables}{\yiHangul}
\setTransitionTo{HangulJamo}{\yiHangul}
\setTransitionTo{HangulJamoExtendedA}{\yiHangul}
\setTransitionTo{HangulJamoExtendedB}{\yiHangul}
\setTransitionTo{HangulCompatibilityJamo}{\yiHangul}
% Han, kana and CJK punctuation: the second hole — this build's face, or the one the file names for CJK.
\setTransitionTo{CJKUnifiedIdeographs}{\yiHan}
\setTransitionTo{CJKUnifiedIdeographsExtensionA}{\yiHan}
\setTransitionTo{CJKUnifiedIdeographsExtensionB}{\yiHan}
\setTransitionTo{CJKUnifiedIdeographsExtensionC}{\yiHan}
\setTransitionTo{CJKUnifiedIdeographsExtensionD}{\yiHan}
\setTransitionTo{CJKUnifiedIdeographsExtensionE}{\yiHan}
\setTransitionTo{CJKUnifiedIdeographsExtensionF}{\yiHan}
\setTransitionTo{CJKUnifiedIdeographsExtensionG}{\yiHan}
\setTransitionTo{CJKUnifiedIdeographsExtensionH}{\yiHan}
\setTransitionTo{CJKCompatibility}{\yiHan}
\setTransitionTo{CJKCompatibilityForms}{\yiHan}
\setTransitionTo{CJKCompatibilityIdeographs}{\yiHan}
\setTransitionTo{CJKCompatibilityIdeographsSupplement}{\yiHan}
\setTransitionTo{CJKRadicalsSupplement}{\yiHan}
\setTransitionTo{CJKStrokes}{\yiHan}
\setTransitionTo{KangxiRadicals}{\yiHan}
\setTransitionTo{CJKSymbolsAndPunctuation}{\yiHan}
\setTransitionTo{IdeographicSymbolsAndPunctuation}{\yiHan}
\setTransitionTo{IdeographicDescriptionCharacters}{\yiHan}
\setTransitionTo{EnclosedCJKLettersAndMonths}{\yiHan}
\setTransitionTo{EnclosedIdeographicSupplement}{\yiHan}
\setTransitionTo{Bopomofo}{\yiHan}
\setTransitionTo{BopomofoExtended}{\yiHan}
\setTransitionTo{Hiragana}{\yiHan}
\setTransitionTo{Katakana}{\yiHan}
\setTransitionTo{KatakanaPhoneticExtensions}{\yiHan}
\setTransitionTo{KanaSupplement}{\yiHan}
\setTransitionTo{KanaExtendedA}{\yiHan}
\setTransitionTo{KanaExtendedB}{\yiHan}
\setTransitionTo{SmallKanaExtension}{\yiHan}
\setTransitionTo{Kanbun}{\yiHan}
\setTransitionTo{HalfwidthAndFullwidthForms}{\yiHan}
\fi
